\section{The general theorem: a precise remaining route}
\label{sec:general-route}

The elementary proof leaves the original general request open. This section
identifies a constructive target that does not demand a height bound, proves
its finite combinatorial last step, and specifies the missing dependencies.
It is a research program, not a claim that its hard theorem has been proved.

\subsection{A short audit of the classical route}

The Subspace Theorem approach uses a Riemann--Roch space of functions with
controlled poles at the boundary. Local estimates turn integral points
into solutions of linear-form inequalities. The Subspace Theorem confines
their images to finitely many proper subspaces; a nonzero function with
bounded poles has only boundedly many zeros, so each subspace contributes
finitely many curve points. Covers and lifting arguments handle the
positive-genus cases with too few boundary points
\cite{CorvajaZannier2002}.

This outline has not verified the strength of those dependencies in iPRA.
Qualitative compactness, arbitrary infinite subsequences, an unlocated
exceptional subspace, or a nonconstructive cover cannot be accepted as
hidden axioms. Each use needs a finite replacement with explicitly coded
data and permitted induction.

Quantitative Subspace Theorems control the \emph{number} of exceptional
subspaces under explicit hypotheses; they do not thereby construct them.
Evertse's Theorems 2.1--2.2 give relevant quantitative forms
\cite{Evertse2010}. The following finite formulation avoids assuming that
a global exceptional set has already been found.

\subsection{A finite covering lemma}

\begin{proposition}[Finite cover to cardinality]
Let $P$ be a decidable predicate on canonical point codes, let $\Phi$ map
point codes to finite vectors over a coded field, and let $M,D\in\N$.
Suppose the following two contracts have been proved.

For every finite list $p_1,\ldots,p_t$ of pairwise distinct points satisfying
$P$, the vectors $\Phi(p_i)$ admit a cover by at most $M$ proper linear
subspaces, with decidable membership in each supplied subspace.

For every subspace $W$ occurring in such a cover, at most $D$ distinct
points in the list have image in $W$.

Then no list of $MD+1$ pairwise distinct points satisfying $P$ exists.
\end{proposition}

\begin{proof}
Suppose $t=MD+1$ distinct points are given. Obtain a cover
$W_1,\ldots,W_s$ with $s\le M$. For each of the finitely many points,
test the subspaces in order and assign it to the least index $j$ for
which $\Phi(p_i)\in W_j$. The covering contract ensures that the finite
search succeeds. No selection from an infinite family occurs.

Let $n_j$ count the points assigned to $j$. Each point is assigned
exactly once, so $t=\sum_{j=1}^s n_j$. Every point counted by $n_j$
lies above $W_j$, and the input points are distinct; hence $n_j\le D$.
Consequently $MD+1=t\le sD\le MD$, a contradiction. These sums,
least indices, and counts are primitive recursive, and their
equalities and inequalities have quantifier-free finite-loop proofs.
\end{proof}

Using the proposition for curves is conditional on proving its two
contracts. It returns no height bound. If there are several local
inequality systems, with respective bounds $M_i,D_i$, and at most $E$
points in an explicitly enumerated exception box, the same assignment
argument gives the cardinality bound
\[
                         E+\sum_i M_iD_i.
\]
No point outside the exception box has thereby been located.

\subsection{Why finite covers can have finite certificates}

Fix a list of $t$ vectors in dimension $n$. If a proper subspace $W$
covers a sublist, replace $W$ by the span of that sublist. This new
space remains inside $W$ and is proper. Finite Gaussian elimination
selects a basis of at most $n-1$ input vectors, including the empty
basis for a zero-dimensional span.

For a particular tuple, therefore, a covering certificate may be sought
among lists of at most $M$ subsets of at most $n-1$ of its $t$ vectors.
There are at most $2^t$ candidate subsets and at most $(2^t+1)^M$
padded lists of candidates. Rank, dependence, and membership have
finite exact checks over rational numbers, or over a supplied
number-field basis. These search limits are primitive recursive.

The existence of a successful certificate for \emph{every} relevant tuple
is still a mathematical theorem. The search can fail; executing it does
not show that it always succeeds. The needed statement is a finite-tuple
version of the quantitative Subspace Theorem, proved with the induction
permitted in iPRA. No such proof is provided here. Assuming a global
finite set of exceptional subspaces would silently introduce stronger
positive information.

\subsection{The geometric input to the fibre bound}

Suppose $\phi_1,\ldots,\phi_n$ are linearly independent rational
functions on a smooth projective curve whose poles are all bounded
by an effective divisor $E$ of degree $D$. A proper linear
subspace admits a nonzero annihilating linear functional. Its
coefficients give a nonzero function $\sum a_i\phi_i$, by linear
independence. Its poles are bounded by $E$. Equality of the degrees
of its zero and pole divisors bounds its distinct zeros by $D$.

This conditional geometric argument suggests the second contract.
It does not construct the functions or their divisor data in arithmetic.
An internal development must represent the function field, divisor,
vector space, and independence certificate, then prove the zero-count
lemma from those finite descriptions. An abstract divisor theorem in
the host Lean library does not automatically discharge that obligation.

\subsection{Explicit unresolved obligations}

A general cardinality proof still needs finite presentations of the
Riemann--Roch functions and their local expansions; a proof of the
finite-tuple quantitative approximation theorem; bounds for its
auxiliary polynomials and matrices using permitted induction; control
of local case splits and the small-height part; and a quantitative
cover-and-lift construction for the positive-genus one- or two-boundary
cases, including all field and denominator changes. Finally, one must
prove the zero-count contract and apply the finite covering lemma.

These are mathematical dependencies. Describing them as standard,
declaring Lean axioms for them, or storing them as unproved structure
fields would not complete the requested proof.

For the positive height/list version, the target is stronger: a
uniform function bounding coordinates in the original affine model,
with a correctness proof. The Runge theorem supplies this for its
explicit family. It does not supply the general bound discussed by
the primary sources. No independence or unprovability conclusion is
drawn from the unresolved state of that step.

